\section{Subdivided trees and abundant extremizers}\label{sec:subdivided}

The sharp bundle bound admits extremizers with arbitrary fixed subcubic
branching patterns. Sufficiently long paths between branching vertices
make the divisor degrees nearly independent of distant branches. An
inequality for joining subtrees then controls subspaces meeting both
sides of a path. Together these arguments prove stability, with no
requirement that the subdivision lengths agree.

For a finite tree $T$ of maximum degree three, suppress its degree-two
vertices; the resulting tree will be called its skeleton. Let $s$ be the
number of trivalent vertices. If $s\ge1$, the skeleton has $s+2$ leaves
and $2s+1$ edges. Replacing its edges by paths of positive integral lengths
gives a subdivision $\Gamma$. Throughout this section put
$b_v=\deg_\Gamma(v)$ in Definition~\ref{graph:construction}.

\begin{theorem}\label{tree:stable}
There are explicit integers $L_s$, depending only on the number $s\ge1$
of trivalent vertices, such that $T_{X(\Gamma)}$ is anticanonically stable
whenever every skeleton edge is replaced by a path of length at least
$L_s$. The lengths may be chosen independently. One can take $L_1=2$
and $L_s=O(\log s)$ as $s\to\infty$.
If $j$ is the total number of edges of $\Gamma$, then
\[
 \dim X(\Gamma)=4j,\qquad \rho(X(\Gamma))=5j+1.
\]
In particular every fixed finite tree of maximum degree three and at
least one edge gives stable tangent bundles after all sufficiently long
independent edge subdivisions.
\end{theorem}

The case with no trivalent vertex is Theorem~\ref{path:stable}.
The examples attain the Picard bound in Theorem~\ref{interact:bound}.
The graph can also be recovered from the fan, so varying its branching
pattern produces distinct extremizers in a common dimension.

\begin{corollary}[Abundance of extremal fans]\label{tree:abundance}
For every integer $q\ge1$, there is an integer $J_q$ such that for every
$j\ge J_q$ there are at least $q$ pairwise nonisomorphic toric fans of
smooth toric Fano varieties with stable tangent bundle, dimension $4j$,
and Picard number $5j+1$. Each belongs to the bundle class of
Theorem~\ref{interact:bound} and attains its maximum in that dimension.
\end{corollary}

\begin{proof}
Choose $q$ finite subcubic trees with an edge and with pairwise distinct
numbers of trivalent vertices, all positive. Such trees exist: take a
path of $s$ vertices and attach leaves so that each path vertex has
degree three, with the three-arm tree for $s=1$. Their branching
patterns remain distinct under subdivision.

For each chosen tree, first suppress degree-two vertices and give every
skeleton edge a length satisfying Theorem~\ref{tree:stable}. Let $j_t$
be the resulting total length for the $t$th tree, and put
$J_q=\max_t j_t$. For any $j\ge J_q$, increase the length of one
skeleton edge in the $t$th tree by $j-j_t$. The stability theorem
allows independent lengths, so every resulting graph bundle is stable.
Each subdivided tree has $j$ edges and $j+1$ vertices, and
$\sum_v b_v=\sum_v\deg(v)=2j$. Proposition~\ref{graph:geometry}
therefore gives dimension $4j$ and Picard number $5j+1$.

Subdivisions of the chosen trees are nonisomorphic because their
numbers of trivalent vertices differ. Lemma~\ref{graph:reconstruction}
then excludes isomorphisms of their toric fans. Finally,
Theorem~\ref{interact:bound} identifies $5j+1$ as the maximum within
the bundle class.
\end{proof}

This corollary concerns isomorphisms of toric fans. It neither classifies
all extremizers nor asserts a growth rate for their number.

\subsection{Subdivision and dependence on distant branches}
The estimate~\eqref{tree:boundary-contraction} makes dependence on
distant branches exponentially small along a subdivided edge.

Convergence of individual ray degrees does not itself imply stability.
The proof below combines it with the rank formula at a cut and an
induction on the number of trivalent vertices. The initial three-arm
case is the quantitative form of Theorem~\ref{spider:stable}. Its
finite rational comparisons, and those for paths, are given in
Appendices~\ref{sec:spiders} and~\ref{sec:paths}; their uniform error
bounds apply to every allowed length. The resulting sufficient lengths
are explicit, though not claimed optimal:
\[
\begin{array}{c|rrrrrrrrrr}
s&1&2&3&4&5&10&25&50&100&1000\\ \hline
L_s&2&27&31&31&33&35&41&43&47&57.
\end{array}
\]
No enumeration of tree shapes is required. This subdivision theorem
concerns the hexagon bundles of Definition~\ref{graph:construction}.
It is distinct from subdivision in the blowup construction, which can
produce instability as in Corollary~\ref{branch:subdivision}.

For clarity, we fix the slice functions used in the proof. Put
$d\nu(t)=(2-|t|)\,dt$ on $[-1,1]$ and
$\mathcal K(s,t)=(3+s-t)^2/2$. The right operator is
\[
 (\mathcal T_Rh)(s)=\int_{-1}^1\mathcal K(s,t)h(t)\,d\nu(t).
\]
Write $h_\infty$ for its positive projective limit, obtained by iterating
from $h_0(t)=2+t$ and normalizing at zero. The contraction estimate
makes the limit independent of the positive initial function. At a
trivalent vertex, integrating the other two branches with incoming
functions $u,v$ gives the outgoing slice function
\[
 J_{u,v}(s)=\iint_{[-1,1]^2}\frac{(4-s-t-z)^3}{6}
                       u(t)v(z)\,d\nu(t)d\nu(z).
\]
This is the cubic integral of \eqref{spider:J}; it is a function of a
boundary coordinate, distinct from the scalar exceptional contributions
$J_{v,e}$ in the blowup construction.

\subsection{Rank at a cut}
Use the independent hexagon reflections and untwisted-base permutations
from Appendix~\ref{sec:spiders}, without permuting branches. An invariant
ray-spanned subspace selects each hexagon as $H_e=0,1,2$ (empty,
distinguished line, or plane), and at each vertex selects the untwisted
rays and twisted ray with indicators $C_v,T_v$. Delete the edges with
$H_e>0$. A component of the remaining forest is called complete for this
selection if $C_v=T_v=1$ at every vertex in it. If $c_0$ counts these
components, the rank is
\begin{equation}\label{tree:rank}
 r=\sum_e H_e+\sum_v b_vC_v+\sum_vT_v-c_0.
\end{equation}
Indeed, first quotient by the selected hexagon directions and untwisted
base coordinates. A selected twisted ray with $C_v=0$ has a private
base coordinate and contributes one dimension. The remaining twisted
rays are rows of the incidence matrix of the zero-edge forest. A
dependence among selected rows extends by zero to a function constant
on each component, and exists precisely on a complete component.

For a subtree rooted at a cut, postpone subtracting the relation of its
root component. Write $a=1$ if this component is complete, and $a=0$
otherwise. Let $e$ record a nonempty selection. Let $f=1$ mean that
every hexagon is a plane, every vertex with $b_v\ge2$ has $C_v=1$,
and every leaf has $C_v\lor T_v=1$. These are precisely the local
requirements for full rank when the tree is completed. Necessity follows
from the independent nontrivial representations at hexagons and bases,
and from the private coordinate of each leaf; they also suffice.
The possible triples are
\[
 (a,e,f)=(0,0,0),(0,1,0),(1,1,0),(0,1,1),(1,1,1).
\]
Let $c_*$ count the complete components other than the root component.
If $w_e(H)$ is the selected hexagon weight and $\delta_v$ one base-ray
weight, define the open deficit
\[
 A=\sum_e(H_e-w_e(H_e))
    +\sum_v(b_vC_v+T_v)(1-\delta_v)-c_*.
\]
This is the rank deficit with the single boundary relation left uncounted.
Joining two rooted pieces with open deficits $A,B$ and indicators $a,b$
across a hexagon $H$ gives final deficit
\begin{equation}\label{tree:join}
 \begin{cases}
 A+B-ab,&H=0,\\
 A+B+H-w(H)-a-b,&H>0.
 \end{cases}
\end{equation}
For $H=0$ the two root components merge and are complete exactly when
$a=b=1$; for $H>0$ they close separately. The combined selection is
nonempty exactly when $e\lor e'\lor(H>0)$, and has full rank exactly
when $f\land f'\land(H=2)$.

\subsection{An inequality inherited from a completed tree}
Consider a finite subtree ending at a degree-two vertex adjacent to the
cut. Give it the limiting degrees obtained by continuing the cut to an
ordinary leaf whose distance tends to infinity. Equivalently, its incoming
function at the cut is $h_\infty$. Let $\alpha$ be its all-selected
open deficit.

\begin{lemma}\label{tree:inherit}
Suppose every sufficiently long completion to that leaf has deficit
at least $\mu>0$ for every invariant selection of rank strictly between
zero and the ambient dimension. For the retained subtree,
\[
\begin{array}{c|c}
\text{boundary condition}&\text{inequality}\\ \hline
a=0,\ e=1&A\ge\mu\\
a=1,\ f=0&A-\alpha\ge\mu\\
a=1,\ f=1&A-\alpha\ge0.
\end{array}
\]
The empty selection has $A=0$.
\end{lemma}

\begin{proof}
For $a=0,e=1$, select no ray in the continuation or its joining
hexagon. The completed selection is proper and nonempty, including
when the retained subtree has $f=1$. Formula~\eqref{tree:join} makes
its deficit exactly $A$. Pass to the limit in the finitely many
retained degrees.

For $a=1,f=0$, select every ray in the continuation and take $H=2$
at the cut. The full-rank failure in the retained subtree persists.
Subtract the all-selected deficit, which is zero. The continuation
and joining terms cancel exactly, leaving $A-\alpha$. Only then
pass to the limit, so no unbounded sum is interchanged with a limit.
For $a=f=1$, adding omitted rays preserves the open rank and boundary
indicator and decreases the deficit, since their weights are positive.
Thus the all-selected configuration minimizes this state.
\end{proof}

\subsection{Nonnegative limiting transitions}
Let $w_1,w_2$ be the limiting weights of a distinguished hexagon line
and full plane, and let $\delta$ be one limiting degree-two-base weight.
Appendix~\ref{sec:paths} gives $w_2+3\delta=4$. The rational moment
matrices at index $40$, with Lemma~\ref{path:contraction}, give
\begin{equation}\label{tree:bulk-bounds}
 \frac{187}{100}<w_2<\frac{15}{8},\qquad 0<w_1<\frac{14}{25}.
\end{equation}
The rational approximations are
\[
 w_1\simeq0.5579289459,\qquad w_2\simeq1.8740976113.
\]
For the exact index-$40$ rational values, with $D=11757/10000$,
$B=81$ and $y=2D(2/7)^{39}$, the errors are
at most $2B(y+y^2)$ and $6B(y+y^2)$, respectively; these bounds
are smaller than the distances to the endpoints in
\eqref{tree:bulk-bounds}.

Put $\beta=w_2/2$ and replace the open deficit by $A-\beta a$.
Extending across one hexagon $H$ and one new degree-two vertex $C,T$
changes the boundary indicator to $a'=CTa$ for $H=0$, and $a'=CT$ for
$H>0$. The adjusted increment is
\begin{equation}\label{tree:increment}
 (2C+T)(1-\delta)+H-w(H)-a\mathbf1_{\{H>0\}}
       +\beta(a-a').
\end{equation}
For $C=T=0$ its three values, indexed by $H=0,1,2$, are
\[
\begin{array}{c|ccc}
a=0&0&1-w_1&2-w_2\\
a=1&\beta&\beta-w_1&1-\beta.
\end{array}
\]
When $CT=0$ and $(C,T)\ne(0,0)$, add
$(2C+T)(1-\delta)$. For $C=T=1$ the values are
\[
\begin{array}{c|ccc}
a=0&w_2-1&\beta-w_1&1-\beta\\
a=1&w_2-1&w_2-w_1-1&0.
\end{array}
\]
All nonzero entries exceed $1/16$ by~\eqref{tree:bulk-bounds}; also
$1-\delta=(w_2-1)/3>0.29$. The only zero transitions are
$(a,C,T,H)=(0,0,0,0)$ and $(1,1,1,2)$, the empty and all-selected
periods. In the same adjusted deficits the joining correction is
\begin{equation}\label{tree:adjusted-join}
 \begin{cases}
 \beta(a+b)-ab,&H=0,\\
 H-w(H)+(\beta-1)(a+b),&H>0.
 \end{cases}
\end{equation}
Its only zeros are $(a,b,H)=(0,0,0)$ and $(1,1,2)$; all other
values exceed $1/16$. Thus arbitrarily many limiting periods may be
inserted without accumulating a negative increment.

\subsection{Uniform control of the geometric error}
Split an internal skeleton edge of length $2q+1+m$, $m\ge0$,
leaving $q$ edges and $q$ degree-two vertices at each end. Call these
retained paths the two necks. The middle consists of $m+1$ hexagons
and $m$ degree-two bases. The two components, capped by a leaf on
the chosen edge, have $s_1,s_2\ge1$ trivalent vertices, with
$s_1+s_2=s$. Retain all their other edge lengths.

Define comparison weights by putting the incoming function $h_\infty$
at each cut and using the limiting weights on the middle. Every
individual ray weight is at most $B=81$, provided all skeleton
lengths are at least two. Indeed a junction slice function
$J_{u,v}$ in~\eqref{spider:J} has oscillation ratio at most $27$.
Each hexagon has at least one endpoint of degree two or one; that
side's function has oscillation at most nine. The facet measure has
mass one and $\nu$ has mass three, giving $27\cdot9/3=81$.
The base-ray bounds are at most $1,2,3$ at leaves, degree-two
vertices and junctions.

Changing one incoming function by Hilbert distance $\eta$ changes
$J_{u,v}$ by at most $\eta$: normalize its pointwise ratio to lie in
$[1,e^\eta]$ and integrate the positive cubic. At each degree-two
vertex the distance contracts by $\tau=2/7$. There is a unique path
from the cut to each vertex. Consequently only one incoming argument
changes at every later junction. The disturbance can enter both
outgoing branches, but does not add to itself at a junction.

Set
\begin{equation}\label{tree:error}
 x_q=D\tau^{q-1},\qquad
 c_s=\frac{156s+272}{5},\qquad
 E_s(q)=B c_s x_q(1+2x_q),
 \qquad D=\frac{11757}{10000}.
\end{equation}

\begin{lemma}\label{tree:comparison}
If all skeleton lengths are at least two and $2x_q\le1$, the
actual and comparison deficits of every selection differ by at most
$E_s(q)$, independently of all subdivision lengths.
\end{lemma}

\begin{proof}
The incoming function at either cut has traversed at least $q$
degree-two vertices from the opposite junction. Its distance from
the limit is at most $x=x_q$. Along the retained neck it contracts
geometrically toward the junction. Six hexagon rays and three base
rays per step give a sum of logarithmic degree errors at most
$9x/(1-\tau)$. The base closest to the cut has error at most $x$;
the adjacent retained hexagon has an additional contraction, which
this estimate discards.

A capped component with $s_i$ junctions has $2s_i+1$ skeleton edges.
One is the neck just counted; each other edge contributes at most
$9x/(1-\tau)$, by contraction along its interior. We discard any
attenuation before reaching that edge. Its $s_i$ junction bases add
$4s_ix$, and its $s_i+1$ genuine leaf bases add $2(s_i+1)x$.
The endpoint total is at most
\[
 \left\{\frac{9(2s_i+1)}{1-\tau}+6s_i+2\right\}x.
\]

Number the chosen bridge edges $i=0,\ldots,L-1$, where
$L=2q+1+m$, and its internal bases $c=1,\ldots,L-1$. On the
middle edges $q\le i\le L-q-1$, the omitted left and right
distances are at most $D\tau^{i-1}$ and $D\tau^{L-2-i}$.
On middle bases $q+1\le c\le L-q-1$, they are at most
$D\tau^{c-2}$ and $D\tau^{L-2-c}$. Each pair is at most $2x$;
summing either set over its positions gives at most $2x/(1-\tau)$.
The six hexagon and three base rays therefore add at most
$18x/(1-\tau)$.

The sum of all logarithmic ray errors is consequently at most
\[
 \left\{\frac{9(2s+4)}{1-\tau}+6s+4\right\}x=c_sx.
\]
Each individual error $\eta$ is at most $2x\le1$, and
$e^\eta-1\le\eta(1+2x)$. Multiplying by $B$ bounds the sum of
absolute ray-weight changes by $E_s(q)$. The selected rank is
unchanged, so the same bound holds for every deficit.
\end{proof}

\subsection{Joining and induction}
Let $\alpha_i$ be the all-selected open deficit of endpoint $i$.
Glue two identical copies of this endpoint with a bridge of length
$2q+1$. Its actual all-selected deficit is zero. The comparison
deficit is $2\alpha_i-w_2$. Lemma~\ref{tree:comparison}, applied
with $2s_i$ junctions and the same $q$, gives
\[
 |\alpha_i-\beta|\le\tfrac12 E_{2s_i}(q).
\]
This uses the degree identity for the doubled tree, without a stability
assumption. Since $c_s$ is affine,
\begin{equation}\label{tree:potential-sum}
 \tfrac12\{E_{2s_1}(q)+E_{2s_2}(q)\}=E_s(q).
\end{equation}

Suppose the two completed endpoint families have invariant margins
$\mu_1,\mu_2\le\mu_0:=2698/100000<1/16$.
Lemma~\ref{tree:inherit} and~\eqref{tree:potential-sum} show that a
positive-margin endpoint contributes its inherited margin minus its
potential error in the adjusted deficit. An endpoint with $a=f=1$
contributes at least minus that error, and an empty endpoint contributes
zero. All intervening increments and the final joining correction
are nonnegative.

If neither endpoint has a positive-margin state, a proper nonempty
global selection must contain a positive increment or joining
correction. Otherwise the zero transitions propagate only empty to
empty or full to full, giving an excluded global selection. Thus the
comparison deficit is at least $\min(\mu_1,\mu_2)-E_s(q)$.
Subtracting the geometric error proves the joining inequality
\begin{equation}\label{tree:gluing-bound}
 \text{actual deficit}\ \ge\min(\mu_1,\mu_2)-2E_s(q).
\end{equation}

\begin{proof}[Proof of Theorem~\ref{tree:stable}]
For $s\ge2$, let $q_s$ be the least positive integer such that
\begin{equation}\label{tree:threshold}
 2x_{q_s}\le1,\qquad
 E_s(q_s)\le\frac{\mu_0}{4s(s-1)},
 \qquad L_s=2q_s+1.
\end{equation}
Put $L_1=2$ and $\mu_s=\frac{\mu_0}{2}(1+1/s)$.
The minimum exists by geometric decay. These thresholds are
nondecreasing: $c_s$ increases and the allowed error decreases.
We prove the stronger assertion that every invariant selection of
rank strictly between zero and $4j$ has deficit at least $\mu_s$.
The case $s=1$ is the quantitative conclusion of
Theorem~\ref{spider:stable}.

For $s\ge2$, remove any internal skeleton edge and cap each exposed
junction by a leaf. The two smaller trees have $s_1,s_2\ge1$
junctions, with $s_1+s_2=s$. Their unchanged lengths satisfy the
inductive thresholds since $L_s\ge L_{s_i}$. Their new leaf edges
may be continued to arbitrary sufficiently large lengths, giving the
families required by Lemma~\ref{tree:inherit}. Retain $q_s$ neck
edges at either end of the chosen original edge; its length is at
least $2q_s+1$. Formula~\eqref{tree:gluing-bound} now gives
\[
 \text{deficit}\ge\mu_{s-1}-2E_s(q_s)
 \ge\mu_{s-1}-\frac{\mu_0}{2s(s-1)}=\mu_s>0.
\]
The independent local symmetries and connected ray matroid therefore
give stability by Proposition~\ref{prop:symmetry}.

For the growth bound, it suffices to choose
\[
 q\ge1+\left\lceil
  \log_{7/2}\!\left(\frac{8BDc_s s(s-1)}{\mu_0}\right)
 \right\rceil.
\]
Then $2x_q\le1$ and
$E_s(q)\le2Bc_sx_q\le\mu_0/[4s(s-1)]$. Thus $L_s=O(\log s)$.
The dimension and Picard number follow from
Proposition~\ref{graph:geometry}, since $\sum_v b_v=2j$ and a tree
has $j+1$ vertices. For a tree with original degree-two vertices,
the skeleton lengths are sums of the original subdivision lengths.
The case $s=0$ is Theorem~\ref{path:stable}.
\end{proof}

The quantitative margin $\mu_s>0.01349$ concerns these locally invariant
selections only. The theorem asserts stability for all subspaces through
the symmetry criterion, without a uniform quantitative claim for them.

\begin{corollary}\label{tree:two-junctions}
Suppose a tree has two trivalent vertices joined by a path of length
at least $27$, with two outer arms at each vertex. If all four outer
arms have length at least two, the graph bundle with $b_v=\deg(v)$ has
anticanonically stable tangent bundle.
\end{corollary}

\begin{proof}
The two completed endpoint families are three-arm trees, so both
inherited margins are $\mu_0$ even when the outer lengths are less
than $L_2$. Lemma~\ref{tree:comparison} needs only outer lengths
at least two. With $q=13$, formula~\eqref{tree:gluing-bound} gives
\[
 \mu_0-2E_2(13)=
 \frac{122113777399878205459586741}
      {5986913480642700450031250000}>0.02039.
\]
\end{proof}
